An agent that trades against rivals has to judge what a deal gives the other side, turn what it acquires into income, and stay solvent. We study these three steps in twenty complete four-player Monopoly games played by four configured LLM systems (6,989 decisions, every game ended by bankruptcy), using legal menus, response records, delivered prompts and canonical events. Bargaining is priced and sensitive to the rival's gain. Across 918 acceptance-legal offers, the side that would complete a color group is the side paying cash in 373 of 429 one-sided cases. Offers that complete only the writer's group are rejected outright in 159 of 289 cases and accepted in 19, while offers that complete only the recipient's group are usually countered. The four systems accept 26 to 30 offers each and differ in how they decline. Completing deeds change hands mainly when their holder is short of cash. In 21 of 48 completing sales the seller held under \$100. A completed group is not yet a developed one. The systems complete 13 to 18 groups each by trade, build on them at the first opportunity in 16 of 17, 17 of 18, 7 of 15 and 3 of 13 cases, and later collect \$20,818, \$26,070, \$5,930 and \$4,560 in rent on them. Mortgaged or illiquid purchases, engine defaults and response-channel conflicts account for 17 of the 20 unbuilt groups. Bankruptcy then moves more property than trade. Fifty-six of sixty eliminations could not have been paid from the debtor's own assets, estates complete 121 groups against 63 by trade, and the eventual winner receives the first estate in 18 games. Seven contextual incidents show what the systems say at these points. The counts describe configured systems under unequal response conditions. They do not rank models, and they do not establish causal effects, optimal play or transfer to other opponents.
